\begin{mistake}{2026-04-19}{03}

\begin{problemcontent}
已知 \(A = \begin{bmatrix} 1 & -2 & 0 \\ 2 & 1 & 5 \\ 0 & 1 & 1 \end{bmatrix}\)，\(A^*\) 是 \(A\) 的伴随矩阵，则 \(A^*x=0\) 的通解是 \underline{\hspace{2cm}}。
\end{problemcontent}

\begin{solutioncontent}
首先计算矩阵 \(A\) 的行列式：
\[
|A| = \begin{vmatrix} 1 & -2 & 0 \\ 2 & 1 & 5 \\ 0 & 1 & 1 \end{vmatrix} = 1 \times (1-5) - (-2) \times (2-0) + 0 = 0
\]
因为 \(A\) 的左上角二阶子式 \(\begin{vmatrix} 1 & -2 \\ 2 & 1 \end{vmatrix} = 5 \neq 0\)，所以 \(r(A) = 2\)。
由伴随矩阵的秩的性质可知，当 \(n=3\) 且 \(r(A) = 2\) 时，\(r(A^*) = 1\)。
因此，齐次线性方程组 \(A^*x=0\) 的基础解系包含 \(3 - r(A^*) = 2\) 个线性无关的解向量。

根据伴随矩阵公式 \(A^*A = |A|E\)，代入 \(|A|=0\) 得 \(A^*A = 0\)。
将 \(A\) 按列分块为 \(A = [\alpha_1, \alpha_2, \alpha_3]\)，则有：
\[
A^* [\alpha_1, \alpha_2, \alpha_3] = [0, 0, 0]
\]
即 \(A^*\alpha_i = 0\ (i=1,2,3)\)，说明 \(A\) 的列向量均是 \(A^*x=0\) 的解。
取 \(A\) 的前两列 \(\alpha_1 = \begin{bmatrix} 1 \\ 2 \\ 0 \end{bmatrix}\) 与 \(\alpha_2 = \begin{bmatrix} -2 \\ 1 \\ 1 \end{bmatrix}\)，二者不成比例，线性无关，可作为基础解系。

故 \(A^*x=0\) 的通解为：
\[
x = c_1 \begin{bmatrix} 1 \\ 2 \\ 0 \end{bmatrix} + c_2 \begin{bmatrix} -2 \\ 1 \\ 1 \end{bmatrix} \quad (c_1, c_2 \text{ 为任意常数})
\]
\end{solutioncontent}

\mistakeknowledge{
\begin{itemize}
 \item \textbf{核心公式/定理}：\(A^*A = |A|E\)；若 \(n\) 阶矩阵 \(r(A)=n-1\)，则 \(r(A^*)=1\)。
 \item \textbf{易错点}：盲目去求伴随矩阵 \(A^*\) 的具体元素，导致计算量过大而出错。
 \item \textbf{关联知识}：矩阵乘积 \(AB=0\) 的列模型视角，即 \(B\) 的每一个列向量都是 \(Ax=0\) 的解。
\end{itemize}}

\end{mistake}
